\SourceLeaf{329}%
\[
\begin{aligned}
l={}&-hr+\tfrac1{24}(2ff'h'+2f''gh'+2fg'h''+2gg''h''\\
&\qquad-ffh'''-2fgh''''-ggh^{\mathrm{v}})r^3\ .\ .\quad(6)
\end{aligned}
\]
Der obigen Bemerkung zufolge darf man nun auch in dieser Gleichung \(l\) mit \(t\) oder mit \(b\) vertauschen. wenn man nur
\[
\begin{array}{rl}
\text{anstatt}&h,\ h',\ h'',\ h''',\ h'''',\ h^{\mathrm{v}}\\
\text{im erstern Falle}&f,\ f',\ f'',\ f''',\ f'''',\ f^{\mathrm{v}}\\
\text{und im andern}&g,\ g',\ g'',\ g''',\ g'''',\ g^{\mathrm{v}}
\end{array}
\]
setzt, so dass man hat
\[
\begin{aligned}
t={}&-fr+\tfrac1{24}(2ff'f'+2f'f''g+2ff''g'+2f''gg''\\
&\qquad-fff'''-2ff''''g-f^{\mathrm{v}}gg)r^3\ .\ .\quad(7)\\
b={}&-gr+\tfrac1{24}(2ff'g'+2f''gg'+2fg'g''+2gg''g''\\
&\qquad-ffg'''-2fgg''''-ggg^{\mathrm{v}})r^3\ .\ .\quad(8)
\end{aligned}
\]

\Article{28}
Die drei Formeln (6), (7), (8) enthalten bereits das Wesentliche zur Auflösung unsrer Aufgabe, so dass zu ihrer Vervollständigung nur noch eine mechanische Rechnung, nemlich die Entwicklung der Werthe der verschiedenen Differentialquotienten und deren Substitution übrig bleibt. Jene Entwicklung gibt, indem wir sofort anstatt der unbestimmten Werthe \(x,x'\) u.\,s.\,w. \(y\) u.\,s.\,w. die zu \(X=T,\ Y=B\) gehörigen bestimmten \(f,f'\) u.\,s.\,w., \(g\) u.\,s.\,w. schreiben, und zur Abkürzung noch setzen
\[
\begin{aligned}
\cos B&=c\\
\sin B&=s\\
\BareRoot{(1-ee\sin B^2)}&=k
\end{aligned}
\]
folgende achtzehn Werthe:
\[
\begin{aligned}
f&=-\frac{k\sin T}{ac}\mathord{.}s\\
f'&=-\frac{k\cos T}{ac}\mathord{.}s\\
f''&=-\frac{\sin T}{akcc}\mathord{.}(1-2eess+ees^4)\\
f'''&=+\frac{k\sin T}{ac}\mathord{.}s\\
f''''&=-\frac{\cos T}{akcc}\mathord{.}(1-2eess+ees^4)\\
f^{\mathrm{v}}&=-\frac{\sin T}{ak^3c^3}\mathord{.}((2-3ee)s+(ee+2e^4)s^3-(2ee+e^4)s^5+e^4s^7)
\end{aligned}
\]
